% This must be in the first 5 lines to tell arXiv to use pdfLaTeX, which is strongly recommended.
\pdfoutput=1
% In particular, the hyperref package requires pdfLaTeX in order to break URLs across lines.

\documentclass[11pt]{article}

% Remove the "review" option to generate the final version.
\usepackage[review]{ACL2023}

% Standard package includes
\usepackage{times}
\usepackage{latexsym}

% For proper rendering and hyphenation of words containing Latin characters (including in bib files)
\usepackage[T1]{fontenc}
% For Vietnamese characters
% \usepackage[T5]{fontenc}
% See https://www.latex-project.org/help/documentation/encguide.pdf for other character sets

% This assumes your files are encoded as UTF8
\usepackage[utf8]{inputenc}

% This is not strictly necessary, and may be commented out.
% However, it will improve the layout of the manuscript,
% and will typically save some space.
\usepackage{microtype}

% This is also not strictly necessary, and may be commented out.
% However, it will improve the aesthetics of text in
% the typewriter font.
\usepackage{inconsolata}


% If the title and author information does not fit in the area allocated, uncomment the following
%
%\setlength\titlebox{<dim>}
%
% and set <dim> to something 5cm or larger.

\title{Instructions for ACL 2023 Proceedings}

% Author information can be set in various styles:
% For several authors from the same institution:
% \author{Author 1 \and ... \and Author n \\
%         Address line \\ ... \\ Address line}
% if the names do not fit well on one line use
%         Author 1 \\ {\bf Author 2} \\ ... \\ {\bf Author n} \\
% For authors from different institutions:
% \author{Author 1 \\ Address line \\  ... \\ Address line
%         \And  ... \And
%         Author n \\ Address line \\ ... \\ Address line}
% To start a seperate ``row'' of authors use \AND, as in
% \author{Author 1 \\ Address line \\  ... \\ Address line
%         \AND
%         Author 2 \\ Address line \\ ... \\ Address line \And
%         Author 3 \\ Address line \\ ... \\ Address line}

\author{First Author \\
  Affiliation / Address line 1 \\
  Affiliation / Address line 2 \\
  Affiliation / Address line 3 \\
  \texttt{email@domain} \\\And
  Second Author \\
  Affiliation / Address line 1 \\
  Affiliation / Address line 2 \\
  Affiliation / Address line 3 \\
  \texttt{email@domain} \\}

\begin{document}
\maketitle
\begin{abstract}
This document is a supplement to the general instructions for *ACL authors. It contains instructions for using the \LaTeX{} style file for ACL 2023.
The document itself conforms to its own specifications, and is, therefore, an example of what your manuscript should look like.
These instructions should be used both for papers submitted for review and for final versions of accepted papers.
\end{abstract}

\section{Introduction}

Large language models (LLMs) are increasingly used for more than answering questions or completing everyday tasks. People also turn to them for advice about personal problems, relationships, and disagreements with others. Research on AI companionship illustrates how users can experience conversational systems as supportive, always available, and nonjudgmental (Pataranutaporn et al., 2025). However, when someone asks an LLM for advice about an interpersonal conflict, being supportive is not always enough. A person describing an argument with a friend, for example, may want reassurance but may also need to hear that their own behavior contributed to the disagreement. A useful response should therefore be able to acknowledge the user's feelings without necessarily approving their actions. This distinction matters because experimental research suggests that overly affirming AI advice can strengthen users' belief that they were right while reducing their willingness to repair interpersonal conflicts (Cheng et al., 2026a).
This tendency to provide excessively agreeable responses is known as sycophancy. It describes situations in which a language model favors a user's beliefs or preferences instead of providing an independent assessment. Sharma et al. (2024) investigated this behavior across several language models and found that their responses were often influenced by users' expressed opinions. The authors also showed that human preference judgments can favor agreeable responses, even when more accurate alternatives are available. However, sycophancy becomes more difficult to recognize when a conversation involves personal experiences or moral judgments rather than factual questions. In an interpersonal disagreement, there may be no single objectively correct answer, and different people may interpret the same behavior differently. Moreover, a model does not need to explicitly agree with the user to favor their perspective. It may do so simply by accepting their interpretation of events or avoiding criticism when criticism would be appropriate.
Cheng et al. (2026b) address this broader problem in ELEPHANT, a benchmark designed to measure social sycophancy in language models. They define social sycophancy as the excessive preservation of a user's desired self-image, or face. This can involve validating the user's perspective, avoiding direct criticism, or presenting questionable behavior in a more favorable way. To investigate these tendencies, ELEPHANT uses different advice-seeking situations, including interpersonal conflicts from Reddit's r/AmITheAsshole (AITA) community. On this forum, people describe a disagreement and ask others to judge whether their behavior was justified. In the AITA-YTA evaluation, the researchers focus on stories where the community judged the narrator to be at fault, giving a "You're the Asshole" (YTA) verdict. Their findings show that LLMs preserve users' face substantially more often than human respondents, even in situations involving questionable behavior. These results demonstrate how models may offer reassuring advice without sufficiently challenging the user's actions. At the same time, Reddit verdicts represent human reference judgments rather than objective moral truth, and expressing empathy should not automatically be interpreted as moral endorsement.
An important question, however, is whether such judgments depend on how a conflict is presented during a conversation. In practice, someone seeking advice may not explain an entire disagreement in a single message. They might first describe why they were upset, then provide additional background, and only later explain what they themselves did. During this exchange, the model responds to each part of the story before knowing the complete situation. These earlier responses then become part of the conversation history. This creates the possibility that an initially sympathetic interpretation influences later judgments, even when new information calls for reconsideration. Alternatively, receiving information gradually may help a model revise its assessment as the situation becomes clearer. Previous research shows that conversational interactions can affect sycophancy. Hong et al. (2025) and Tang et al. (2026) investigate how models respond to repeated user pressure, while Jain et al. (2026) examine the influence of previous interaction context. More closely related, Wu et al. (2026) compare single-turn and five-turn presentations of interpersonal conflicts and find greater alignment with the narrator following progressive disclosure. These findings suggest that conversational structure matters, but they leave room to investigate how moral verdicts vary across additional dialogue lengths and across conflicts with different human reference judgments. In particular, it remains relevant to examine whether extending a conversation from one to five or ten turns changes the model's final assessment when the underlying conflict remains the same.
Another relevant question concerns the relationship between a model's generated reasoning and its final answer. A model might express doubts about the user's behavior in its reasoning while still giving a reassuring final verdict, or it might construct reasoning that consistently supports the user. Turpin et al. (2023) demonstrated that chain-of-thought explanations do not always faithfully reflect the factors influencing a model's answers. More recently, Feng et al. (2026) found that reasoning can reduce sycophancy in final decisions but may also produce biased or one-sided justifications. These findings motivate examining whether the moral judgments expressed in generated reasoning are consistent with final answers, and whether such patterns differ across language models. Importantly, generated reasoning should not be interpreted as direct evidence of a model's internal decision-making process.
Building on this research, we investigate how presenting the same interpersonal conflict across different numbers of conversational turns affects an LLM's tendency to produce user-favoring moral judgments. Our study uses 500 AITA stories derived from ELEPHANT, balanced between 250 cases labeled "You're the Asshole" (YTA) and 250 labeled "Not the Asshole" (NTA) by the Reddit community. We compare three conversational conditions: the complete story presented in one prompt, or the same conflict revealed progressively over five or ten user turns. In the multi-turn conditions, the model responds between parts of the story before receiving the final request for a moral verdict. We examine particularly whether models produce NTA judgments for stories that received YTA community verdicts, while using NTA-labeled stories to account for general tendencies to favor the narrator. Such disagreements are treated as indicators of potentially user-favoring judgments rather than definitive evidence of sycophancy. We additionally explore reasoning–answer consistency and differences across models where the available outputs permit these analyses. Through these comparisons, our study aims to better understand whether, and in what direction, the way an interpersonal conflict unfolds over a conversation influences the moral judgments produced by language models.

\section{Engines}

To produce a PDF file, pdf\LaTeX{} is strongly recommended (over original \LaTeX{} plus dvips+ps2pdf or dvipdf). Xe\LaTeX{} also produces PDF files, and is especially suitable for text in non-Latin scripts.
\begin{table}
\centering
\begin{tabular}{lc}
\hline
\textbf{Command} & \textbf{Output}\\
\hline
\verb|{\"a}| & {\"a} \\
\verb|{\^e}| & {\^e} \\
\verb|{\`i}| & {\`i} \\ 
\verb|{\.I}| & {\.I} \\ 
\verb|{\o}| & {\o} \\
\verb|{\'u}| & {\'u}  \\ 
\verb|{\aa}| & {\aa}  \\\hline
\end{tabular}
\begin{tabular}{lc}
\hline
\textbf{Command} & \textbf{Output}\\
\hline
\verb|{\c c}| & {\c c} \\ 
\verb|{\u g}| & {\u g} \\ 
\verb|{\l}| & {\l} \\ 
\verb|{\~n}| & {\~n} \\ 
\verb|{\H o}| & {\H o} \\ 
\verb|{\v r}| & {\v r} \\ 
\verb|{\ss}| & {\ss} \\
\hline
\end{tabular}
\caption{Example commands for accented characters, to be used in, \emph{e.g.}, Bib\TeX{} entries.}
\label{tab:accents}
\end{table}
\section{Preamble}
\begin{table*}
\centering
\begin{tabular}{lll}
\hline
\textbf{Output} & \textbf{natbib command} & \textbf{Old ACL-style command}\\
\hline
\citep{ct1965} & \verb|\citep| & \verb|\cite| \\
\citealp{ct1965} & \verb|\citealp| & no equivalent \\
\citet{ct1965} & \verb|\citet| & \verb|\newcite| \\
\citeyearpar{ct1965} & \verb|\citeyearpar| & \verb|\shortcite| \\
\citeposs{ct1965} & \verb|\citeposs| & no equivalent \\
\citep[FFT;][]{ct1965} &  \verb|\citep[FFT;][]| & no equivalent\\
\hline
\end{tabular}
\caption{\label{citation-guide}
Citation commands supported by the style file.
The style is based on the natbib package and supports all natbib citation commands.
It also supports commands defined in previous ACL style files for compatibility.
}
\end{table*}
The first line of the file must be
\begin{quote}
\begin{verbatim}
\documentclass[11pt]{article}
\end{verbatim}
\end{quote}
To load the style file in the review version:
\begin{quote}
\begin{verbatim}
\usepackage[review]{ACL2023}
\end{verbatim}
\end{quote}
For the final version, omit the \verb|review| option:
\begin{quote}
\begin{verbatim}
\usepackage{ACL2023}
\end{verbatim}
\end{quote}
To use Times Roman, put the following in the preamble:
\begin{quote}
\begin{verbatim}
\usepackage{times}
\end{verbatim}
\end{quote}
(Alternatives like txfonts or newtx are also acceptable.)
Please see the \LaTeX{} source of this document for comments on other packages that may be useful.
Set the title and author using \verb|\title| and \verb|\author|. Within the author list, format multiple authors using \verb|\and| and \verb|\And| and \verb|\AND|; please see the \LaTeX{} source for examples.
By default, the box containing the title and author names is set to the minimum of 5 cm. If you need more space, include the following in the preamble:
\begin{quote}
\begin{verbatim}
\setlength\titlebox{<dim>}
\end{verbatim}
\end{quote}
where \verb|<dim>| is replaced with a length. Do not set this length smaller than 5 cm.

\section{Document Body}

\subsection{Footnotes}

Footnotes are inserted with the \verb|\footnote| command.\footnote{This is a footnote.}

\subsection{Tables and figures}

See Table~\ref{tab:accents} for an example of a table and its caption.
\textbf{Do not override the default caption sizes.}

\subsection{Hyperlinks}

Users of older versions of \LaTeX{} may encounter the following error during compilation: 
\begin{quote}
\tt\verb|\pdfendlink| ended up in different nesting level than \verb|\pdfstartlink|.
\end{quote}
This happens when pdf\LaTeX{} is used and a citation splits across a page boundary. The best way to fix this is to upgrade \LaTeX{} to 2018-12-01 or later.

\subsection{Citations}



Table~\ref{citation-guide} shows the syntax supported by the style files.
We encourage you to use the natbib styles.
You can use the command \verb|\citet| (cite in text) to get ``author (year)'' citations, like this citation to a paper by \citet{Gusfield:97}.
You can use the command \verb|\citep| (cite in parentheses) to get ``(author, year)'' citations \citep{Gusfield:97}.
You can use the command \verb|\citealp| (alternative cite without parentheses) to get ``author, year'' citations, which is useful for using citations within parentheses (e.g. \citealp{Gusfield:97}).

\subsection{References}

\nocite{Ando2005,augenstein-etal-2016-stance,andrew2007scalable,rasooli-tetrault-2015,goodman-etal-2016-noise,harper-2014-learning}

The \LaTeX{} and Bib\TeX{} style files provided roughly follow the American Psychological Association format.
If your own bib file is named \texttt{custom.bib}, then placing the following before any appendices in your \LaTeX{} file will generate the references section for you:
\begin{quote}
\begin{verbatim}
\bibliographystyle{acl_natbib}
\bibliography{custom}
\end{verbatim}
\end{quote}
You can obtain the complete ACL Anthology as a Bib\TeX{} file from \url{https://aclweb.org/anthology/anthology.bib.gz}.
To include both the Anthology and your own .bib file, use the following instead of the above.
\begin{quote}
\begin{verbatim}
\bibliographystyle{acl_natbib}
\bibliography{anthology,custom}
\end{verbatim}
\end{quote}
Please see Section~\ref{sec:bibtex} for information on preparing Bib\TeX{} files.

\subsection{Appendices}

Use \verb|\appendix| before any appendix section to switch the section numbering over to letters. See Appendix~\ref{sec:appendix} for an example.

\section{Bib\TeX{} Files}
\label{sec:bibtex}

Unicode cannot be used in Bib\TeX{} entries, and some ways of typing special characters can disrupt Bib\TeX's alphabetization. The recommended way of typing special characters is shown in Table~\ref{tab:accents}.

Please ensure that Bib\TeX{} records contain DOIs or URLs when possible, and for all the ACL materials that you reference.
Use the \verb|doi| field for DOIs and the \verb|url| field for URLs.
If a Bib\TeX{} entry has a URL or DOI field, the paper title in the references section will appear as a hyperlink to the paper, using the hyperref \LaTeX{} package.

\section*{Limitations}
ACL 2023 requires all submissions to have a section titled ``Limitations'', for discussing the limitations of the paper as a complement to the discussion of strengths in the main text. This section should occur after the conclusion, but before the references. It will not count towards the page limit.
The discussion of limitations is mandatory. Papers without a limitation section will be desk-rejected without review.

While we are open to different types of limitations, just mentioning that a set of results have been shown for English only probably does not reflect what we expect. 
Mentioning that the method works mostly for languages with limited morphology, like English, is a much better alternative.
In addition, limitations such as low scalability to long text, the requirement of large GPU resources, or other things that inspire crucial further investigation are welcome.

\section*{Ethics Statement}
Scientific work published at ACL 2023 must comply with the ACL Ethics Policy.\footnote{\url{https://www.aclweb.org/portal/content/acl-code-ethics}} We encourage all authors to include an explicit ethics statement on the broader impact of the work, or other ethical considerations after the conclusion but before the references. The ethics statement will not count toward the page limit (8 pages for long, 4 pages for short papers).

\section*{Acknowledgements}
This document has been adapted by Jordan Boyd-Graber, Naoaki Okazaki, Anna Rogers from the style files used for earlier ACL, EMNLP and NAACL proceedings, including those for
EACL 2023 by Isabelle Augenstein and Andreas Vlachos,
EMNLP 2022 by Yue Zhang, Ryan Cotterell and Lea Frermann,
ACL 2020 by Steven Bethard, Ryan Cotterell and Rui Yan,
ACL 2019 by Douwe Kiela and Ivan Vuli\'{c},
NAACL 2019 by Stephanie Lukin and Alla Roskovskaya, 
ACL 2018 by Shay Cohen, Kevin Gimpel, and Wei Lu, 
NAACL 2018 by Margaret Mitchell and Stephanie Lukin,
Bib\TeX{} suggestions for (NA)ACL 2017/2018 from Jason Eisner,
ACL 2017 by Dan Gildea and Min-Yen Kan, NAACL 2017 by Margaret Mitchell, 
ACL 2012 by Maggie Li and Michael White, 
ACL 2010 by Jing-Shin Chang and Philipp Koehn, 
ACL 2008 by Johanna D. Moore, Simone Teufel, James Allan, and Sadaoki Furui, 
ACL 2005 by Hwee Tou Ng and Kemal Oflazer, 
ACL 2002 by Eugene Charniak and Dekang Lin, 
and earlier ACL and EACL formats written by several people, including
John Chen, Henry S. Thompson and Donald Walker.
Additional elements were taken from the formatting instructions of the \emph{International Joint Conference on Artificial Intelligence} and the \emph{Conference on Computer Vision and Pattern Recognition}.

% Entries for the entire Anthology, followed by custom entries
\bibliography{anthology,custom}
\bibliographystyle{acl_natbib}

\appendix

\section{Example Appendix}
\label{sec:appendix}

This is a section in the appendix.

\end{document}
